\documentclass[a4paper,twoside]{article}

\usepackage{morewrites}

\usepackage[svgnames,dvipsnames,x11names]{xcolor}
\usepackage{graphicx}
\usepackage{texsmith-lists}
\usepackage[english]{babel}
\usepackage[colorlinks=true, linkcolor=NavyBlue, urlcolor=NavyBlue, citecolor=NavyBlue]{hyperref}
\usepackage[a4paper]{geometry}
\usepackage{ts-fonts}
\usepackage{booktabs}
\usepackage{csquotes}
\usepackage{tabularx}
\usepackage{amsmath}
\usepackage{float}
\usepackage{seqsplit}
\newcolumntype{Y}{>{\centering\arraybackslash}X}
\newcolumntype{Z}{>{\raggedleft\arraybackslash}X}
\usepackage{ts-callouts}
\usepackage{ts-code}

\usepackage{lastpage}
\usepackage{refcount}
\makeatletter
\def\ps@plain{
\let\@oddhead\@empty
\let\@evenhead\@empty
\def\@oddfoot{
\hfil
\ifnum\getpagerefnumber{LastPage}>1\relax
\thepage
\fi
\hfil}
\let\@evenfoot\@oddfoot
}
\makeatother

\title{Cross-document references}
\author{}\date{}
\begin{document}
\maketitle
A counter number only exists inside the conversion that allocated it. A second
document citing \tscodeinline{FW-\allowbreak{}10} has no way to know what \tscodeinline{FW-\allowbreak{}10} is, and --- worse --- no way
to notice when a renumbering turns it into \tscodeinline{FW-\allowbreak{}12}. Hard-coded numbers in a
sister document break silently, which is exactly the maintenance
custom counters remove \emph{inside} a document.

TeXSmith closes the loop the way DocBook's \emph{target database} and Sphinx's
\tscodeinline{objects.inv} do: every conversion publishes a small JSON \textbf{inventory} next to
its output, and a citing document declares the inventories it depends on.
\begin{tscode}[lang=yaml, engine=pygments]
\begin{Verbatim}[commandchars=\\\{\},breaklines, breakanywhere, commandchars=\\\{\}]
\PY{c+c1}{\PYZsh{} firmware\PYZhy{}review.md — the document that publishes}
\PY{n+nn}{\PYZhy{}\PYZhy{}\PYZhy{}}
\PY{n+nt}{title}\PY{p}{:}\PY{+w}{ }\PY{l+lScalar+lScalarPlain}{Revue}\PY{l+lScalar+lScalarPlain}{ }\PY{l+lScalar+lScalarPlain}{firmware}
\PY{n+nt}{id}\PY{p}{:}\PY{+w}{ }\PY{l+lScalar+lScalarPlain}{RHE\PYZhy{}423}
\PY{n+nt}{press}\PY{p}{:}
\PY{+w}{  }\PY{n+nt}{declare}\PY{p}{:}
\PY{+w}{    }\PY{n+nt}{counters}\PY{p}{:}
\PY{+w}{      }\PY{n+nt}{fw}\PY{p}{:}
\PY{+w}{        }\PY{n+nt}{name}\PY{p}{:}\PY{+w}{ }\PY{l+lScalar+lScalarPlain}{Constat}
\PY{+w}{        }\PY{n+nt}{format}\PY{p}{:}\PY{+w}{ }\PY{l+s}{\PYZdq{}}\PY{l+s}{FW\PYZhy{}\PYZob{}n:02d\PYZcb{}}\PY{l+s}{\PYZdq{}}
\PY{n+nn}{\PYZhy{}\PYZhy{}\PYZhy{}}
\end{Verbatim}
\end{tscode}
\begin{tscode}[lang=yaml, engine=pygments]
\begin{Verbatim}[commandchars=\\\{\},breaklines, breakanywhere, commandchars=\\\{\}]
\PY{c+c1}{\PYZsh{} hardware\PYZhy{}review.md — the document that cites}
\PY{n+nn}{\PYZhy{}\PYZhy{}\PYZhy{}}
\PY{n+nt}{title}\PY{p}{:}\PY{+w}{ }\PY{l+lScalar+lScalarPlain}{Revue}\PY{l+lScalar+lScalarPlain}{ }\PY{l+lScalar+lScalarPlain}{hardware}
\PY{n+nt}{press}\PY{p}{:}
\PY{+w}{  }\PY{n+nt}{sources}\PY{p}{:}
\PY{+w}{    }\PY{n+nt}{crossrefs}\PY{p}{:}
\PY{+w}{      }\PY{n+nt}{fwrev}\PY{p}{:}\PY{+w}{ }\PY{l+lScalar+lScalarPlain}{build/firmware\PYZhy{}review.refs.json}
\PY{n+nn}{\PYZhy{}\PYZhy{}\PYZhy{}}
\end{Verbatim}
\end{tscode}
\begin{tscode}[lang=md, engine=pygments]
\begin{Verbatim}[commandchars=\\\{\},breaklines, breakanywhere, commandchars=\\\{\}]
L\PYZsq{}écart est écrasé par l\PYZsq{}étalement spectral (voir \PY{n+ni}{@fwrev:fw:pas\PYZhy{}de\PYZhy{}temps}).
\end{Verbatim}
\end{tscode}
renders as
\begin{displayquote}
L'écart est écrasé par l'étalement spectral (voir RHE-423-FW-10 p. 14).
\end{displayquote}
\section{Identifying a document}
\tscodeinline{id} is a free label --- a contract number, a report reference, whatever your
organisation numbers documents with (\tscodeinline{document-\allowbreak{}id} is accepted as an alias). When the target document declares one,
a citation concatenates it with the item's label: \tscodeinline{RHE-\allowbreak{}423} + \tscodeinline{FW-\allowbreak{}10} \tsscript{symbols}{\(\rightarrow\)}
\textbf{\tscodeinline{RHE-\allowbreak{}423-\allowbreak{}FW-\allowbreak{}10}}. That is what makes the reference unambiguous outside the
document that defines it.

It is optional. Without it, a citation falls back to naming the document by its
title --- \tscodeinline{FW-\allowbreak{}10 (Revue firmware, p. 14)} --- and without a title either, to the
bare label.
\section{The inventory}
Each conversion writes \tscodeinline{<document>.refs.json} beside its \tscodeinline{.tex}:
\begin{tscode}[lang=json, engine=pygments]
\begin{Verbatim}[commandchars=\\\{\},breaklines, breakanywhere, commandchars=\\\{\}]
\PY{p}{\PYZob{}}
\PY{+w}{  }\PY{n+nt}{\PYZdq{}schema\PYZdq{}}\PY{p}{:}\PY{+w}{ }\PY{l+m+mi}{1}\PY{p}{,}
\PY{+w}{  }\PY{n+nt}{\PYZdq{}document\PYZdq{}}\PY{p}{:}\PY{+w}{ }\PY{p}{\PYZob{}}
\PY{+w}{    }\PY{n+nt}{\PYZdq{}id\PYZdq{}}\PY{p}{:}\PY{+w}{ }\PY{l+s+s2}{\PYZdq{}RHE\PYZhy{}423\PYZdq{}}\PY{p}{,}
\PY{+w}{    }\PY{n+nt}{\PYZdq{}title\PYZdq{}}\PY{p}{:}\PY{+w}{ }\PY{l+s+s2}{\PYZdq{}Revue firmware\PYZdq{}}\PY{p}{,}
\PY{+w}{    }\PY{n+nt}{\PYZdq{}output\PYZdq{}}\PY{p}{:}\PY{+w}{ }\PY{l+s+s2}{\PYZdq{}firmware\PYZhy{}review.pdf\PYZdq{}}\PY{p}{,}
\PY{+w}{    }\PY{n+nt}{\PYZdq{}source\PYZdq{}}\PY{p}{:}\PY{+w}{ }\PY{l+s+s2}{\PYZdq{}../firmware\PYZhy{}review.md\PYZdq{}}\PY{p}{,}
\PY{+w}{    }\PY{n+nt}{\PYZdq{}source\PYZus{}sha256\PYZdq{}}\PY{p}{:}\PY{+w}{ }\PY{l+s+s2}{\PYZdq{}8b66c613…\PYZdq{}}
\PY{+w}{  }\PY{p}{\PYZcb{},}
\PY{+w}{  }\PY{n+nt}{\PYZdq{}anchors\PYZdq{}}\PY{p}{:}\PY{+w}{ }\PY{p}{\PYZob{}}
\PY{+w}{    }\PY{n+nt}{\PYZdq{}fw:pas\PYZhy{}de\PYZhy{}temps\PYZdq{}}\PY{p}{:}\PY{+w}{ }\PY{p}{\PYZob{}}\PY{+w}{ }\PY{n+nt}{\PYZdq{}counter\PYZdq{}}\PY{p}{:}\PY{+w}{ }\PY{l+s+s2}{\PYZdq{}fw\PYZdq{}}\PY{p}{,}\PY{+w}{ }\PY{n+nt}{\PYZdq{}label\PYZdq{}}\PY{p}{:}\PY{+w}{ }\PY{l+s+s2}{\PYZdq{}FW\PYZhy{}10\PYZdq{}}\PY{p}{,}\PY{+w}{ }\PY{n+nt}{\PYZdq{}page\PYZdq{}}\PY{p}{:}\PY{+w}{ }\PY{l+m+mi}{14}\PY{+w}{ }\PY{p}{\PYZcb{}}
\PY{+w}{  }\PY{p}{\PYZcb{}}
\PY{p}{\PYZcb{}}
\end{Verbatim}
\end{tscode}
It is delivered next to the artifact you asked for: in the directory for
\tscodeinline{-\allowbreak{}o dir}, beside the PDF for \tscodeinline{-\allowbreak{}o report.pdf}. Its \tscodeinline{document.source} is stored
relative to its own location and is rewritten when it moves, so the staleness
check keeps working.

It is written in two passes, because the two halves of a reference become known
at different times. The keys and their labels are known while converting; the
\textbf{page numbers only exist once \LaTeX{} has run}, and are harvested from the
\tscodeinline{.aux} afterwards --- the same source the \tscodeinline{xr} package reads. A conversion without
\tscodeinline{-\allowbreak{}-\allowbreak{}build} therefore publishes an inventory without pages, and citations simply
omit the page.
\begin{tscallout}[kind=tip, title={Commit the inventory}]
It is a build artifact, but committing it makes renumbering \emph{diffable}:
\tscodeinline{git diff} shows \tscodeinline{FW-\allowbreak{}04 \(\rightarrow\) FW-\allowbreak{}10} before you publish. For a contractual
document, that diff is the review gate that catches a renumbering before it
reaches a reader.
\end{tscallout}
\section{Declaring and citing}
\begin{tscode}[lang=yaml, engine=pygments]
\begin{Verbatim}[commandchars=\\\{\},breaklines, breakanywhere, commandchars=\\\{\}]
\PY{n+nt}{press}\PY{p}{:}
\PY{+w}{  }\PY{n+nt}{sources}\PY{p}{:}
\PY{+w}{    }\PY{n+nt}{crossrefs}\PY{p}{:}
\PY{+w}{      }\PY{n+nt}{fwrev}\PY{p}{:}\PY{+w}{ }\PY{l+lScalar+lScalarPlain}{build/firmware\PYZhy{}review.refs.json}\PY{+w}{ }\PY{c+c1}{\PYZsh{} shorthand}
\PY{+w}{      }\PY{n+nt}{hwrev}\PY{p}{:}
\PY{+w}{        }\PY{n+nt}{inventory}\PY{p}{:}\PY{+w}{ }\PY{l+lScalar+lScalarPlain}{../hardware/build/hardware\PYZhy{}review.refs.json}
\end{Verbatim}
\end{tscode}
A top-level \tscodeinline{crossrefs:} key is the deprecated spelling; \tscodeinline{tmark lint -\allowbreak{}-\allowbreak{}fix}
moves it under \tscodeinline{press.sources}.

Paths are relative to the citing document. The alias is then the first segment
of the reference: \tscodeinline{@fwrev:fw:pas-\allowbreak{}de-\allowbreak{}temps}.

Resolution is \textbf{explicit}: an alias never falls back to a local counter, and a
local counter can never be shadowed by an external one. \tscodeinline{@fw:x} is always this
document's item, \tscodeinline{@fwrev:fw:x} is always the other document's.
\section{An external citation is text, not a link}
The target lives in a different PDF, so TeXSmith emits the formatted reference
as plain text rather than a \tscodeinline{\textbackslash{}hyperref} that would resolve to nothing. Local
references keep their live link.
\section{Diagnostics}
This is the part that earns the feature. All of these are
diagnostics, shown by default and promotable to hard
failures with \tscodeinline{-\allowbreak{}-\allowbreak{}strict} (or \tscodeinline{press.features.strict: true}):
\begin{center}
\begin{tabularx}{\linewidth}{>{\raggedright\arraybackslash}X>{\raggedright\arraybackslash}X}
\toprule
\textbf{Situation} & \textbf{Behaviour} \\
\midrule
the key is not published by the declared inventory & warning, and the citation renders as \tscodeinline{[?fwrev:fw:disparu]} \\
the inventory file does not exist yet & warning, citations render as \tscodeinline{[?…]} \\
the inventory's source changed since it was written & warning: the numbers you are citing may already be wrong \\
the inventory's source no longer resolves at all & warning: staleness can no longer be checked \\
the inventory uses a newer schema & warning, inventory ignored \\
\bottomrule
\end{tabularx}
\end{center}
Warnings are attributed to the \textbf{citing document}, not to a TeXSmith source
line:
\begin{tscode}[lang=text, engine=pygments]
\begin{Verbatim}[commandchars=\\\{\},breaklines, breakanywhere, commandchars=\\\{\}]
review/hardware\PYZhy{}review.md: warning ref\PYZhy{}unresolved: Cross\PYZhy{}reference \PYZsq{}@fwrev:fw:disparu\PYZsq{} is not published by …
\end{Verbatim}
\end{tscode}
An unresolved citation is deliberately \textbf{visible in the output}. A silently
dropped reference is how a contractual document ships a hole.
\section{Build order}
\tscodeinline{hardware-\allowbreak{}review.pdf} now depends on \tscodeinline{firmware-\allowbreak{}review.refs.json}, which is a
Makefile edge:
\begin{tscode}[lang=makefile, engine=pygments]
\begin{Verbatim}[commandchars=\\\{\},breaklines, breakanywhere, commandchars=\\\{\}]
\PY{n+nf}{build/hardware\PYZhy{}review.pdf}\PY{o}{:}\PY{+w}{ }\PY{n}{hardware}\PYZhy{}\PY{n}{review}.\PY{n}{md} \PY{n}{build}/\PY{n}{firmware}\PYZhy{}\PY{n}{review}.\PY{n}{refs}.\PY{n}{json}
\end{Verbatim}
\end{tscode}
Two documents citing each other need two passes, exactly like \LaTeX{}'s own
\tscodeinline{.aux}: a missing inventory is a warning and not a failure, so the first build
produces the inventories and the second resolves the citations.
\section{Not implemented yet}
Links from one PDF into another (\tscodeinline{xr-\allowbreak{}hyper} does this with a PDF-viewer \tscodeinline{GoToR}
action), sections, figures and tables in the inventory (only counters are
published yet), custom citation formats, and any form of automatic build
ordering.
\end{document}
